% Generated draft: requires mathematical and source review.
\documentclass[11pt]{article}
\usepackage[margin=25mm]{geometry}
\usepackage{fontspec}
\setmainfont{Latin Modern Roman}
\usepackage{amsmath,amssymb,mathtools,unicode-math}
\setmathfont{Latin Modern Math}
\usepackage{enumitem,needspace,xurl,hyperref}
\setlength{\parindent}{0pt}
\setlength{\parskip}{5pt}
\setlength{\emergencystretch}{3em}
\newcommand{\N}{\mathbb{N}}
\newcommand{\Z}{\mathbb{Z}}
\newcommand{\Q}{\mathbb{Q}}
\newcommand{\R}{\mathbb{R}}
\newcommand{\C}{\mathbb{C}}
\newcommand{\F}{\mathbb{F}}
\newcommand{\A}{\mathbb{A}}
\newcommand{\PP}{\mathbb{P}}
\newcommand{\E}{\mathbb{E}}
\newcommand{\eps}{\varepsilon}
\newcommand{\dd}{\,\mathrm{d}}
\begin{document}
\section*{UnsolvedMath ID: 30000961}
\textbf{Source problem code:} catalogue entry
\subsection{Definitions and mathematical statement}
Let $K\subset S^3$ be a knot and let $G$ be a complex reductive group with Lie algebra $\mathfrak g$. Use a compact knot exterior $M=S^3\setminus\operatorname{int}\nu(K)$, with boundary torus $T=\partial M$; its fundamental group agrees with that of $S^3\setminus K$. In the affine GIT convention put $X_G(N)=\operatorname{Hom}(\pi_1(N),G)//G$, and let \[Z_{G,K}=\overline{\operatorname{im}(X_G(M)\to X_G(T))}^{\mathrm{Zar}}.\] Write $A_{G,K}$ for the vanishing ideal of this closed image in the intended boundary-character coordinate ring. Thus $A_{G,K}$ is an ideal, not an unspecified single $A$-polynomial. The clean/current question uses the report's Weyl-invariant recursive ideal $I^W_{\mathfrak g,K}$ and its specialization $\varepsilon(I^W_{\mathfrak g,K})$ to ask whether \[V\bigl(\varepsilon(I^W_{\mathfrak g,K})\bigr)=Z_{G,K},\qquad\text{equivalently}\qquad\sqrt{\varepsilon(I^W_{\mathfrak g,K})}=A_{G,K}.\] Here $\sqrt J=\{a:a^m\in J\text{ for some }m\geq1\}$. The equivalence is the usual vanishing-set/radical correspondence provided that the specialization really defines an ideal in the same complex affine coordinate ring. The extract does not define the quantum invariant sequence, the recursive-operator algebra and ideal, the Weyl action, or the specialization map. A usual classical specialization sends a quantum parameter to one, but its integral form and precise target cannot be inferred from the symbol $\varepsilon$ alone. Nor is a maximal-torus quotient silently identified with the full character variety for every reductive $G$; the relevant components and identifications must be specified.\par The original asks to make explicit a relation labeled Conjecture 7, building on $q$-holonomicity, whereas the verification note attributes the restored radical identity to Conjecture 12. Both labels are retained without reconciliation. The supplied 2026-08-21 triage labels the record open and describes the explicit quantum/character-variety relation as an unresolved general research program rather than a theorem.
\subsection{Short English statement}
Ask whether the specialized Weyl-invariant recursion ideal cuts out exactly the closed boundary-character image of a knot exterior, equivalently whether its radical is the image's vanishing ideal. The recursion algebra, specialization, and common coordinate-ring identification must be supplied.
\subsection{Sources}
\begin{itemize}\item Local UnsolvedMath registry, permanent ID 30000961, source problem catalogue entry, \emph{Recursive Ideals and Knot Character Varieties}, proposed year 2008; no individual proposer supplied. The original refers to Conjecture 7, while the verification note refers to Conjecture 12.\item Supplied citation: \emph{Invariants in Low-Dimensional Topology} (2009), Oberwolfach Reports, DOI 10.4171/owr/2008/22, pp.\ 1157--1232: \url{https://doi.org/10.4171/owr/2008/22}. Supplied literature label: original Oberwolfach report.\item Supplied literature label: Sikora 2008, \url{https://arxiv.org/abs/0807.0943}.\end{itemize}
\end{document}
